\subsection{The norm comparison over every finite coefficient field}
\label{sec:kt-finite-coefficient-norm}

The scalar restriction in the cyclic norm construction can be
removed by a finite polarization identity.  No extension with
additive Galois group \(k\) is treated as cyclic.

\begin{proposition}[KT-CUBIC-FINITE-COEFFICIENT-NORM]
Retain the hypotheses of KT-ORDINARY-CUBIC-MASSEY-CONDITION.
Over \(\mathbf F_p\), put
\[
\begin{split}
 U&=H^1(G_{\mathbf Q,S},\mathbf F_p(\eta\psi)),\\
 V&=\ker\left(H^1(G_{\mathbf Q,S},\mathbf F_p(\psi^{-1}))
       \longrightarrow H^1(\mathbf Q_p,\mathbf F_p(\psi^{-1}))
           \oplus\bigoplus_{q\mid D}H^1(I_q,\mathbf F_p(\psi^{-1}))
         \right).
\end{split}                                               \tag{FC.1}
\]
These are finite-dimensional spaces.  Choose bases
\((u_i)_i\) of \(U\) and \((v_j)_j\) of \(V\).
For any finite \(k/\mathbf F_p\), write
\(u=\sum_i c_i u_i\), \(v=\sum_j d_jv_j\)
under scalar extension.  The actual relative obstruction
\(q(u,v)=\mathfrak o_3(u,v)\) satisfies
\[
\begin{split}
q(u,v)={}&\sum_{i,j}c_i^2d_jq(u_i,v_j)\\
 &+\sum_{i<i',j}c_ic_{i'}d_j
     \bigl(q(u_i+u_{i'},v_j)-q(u_i,v_j)-q(u_{i'},v_j)\bigr).
\end{split}                                               \tag{FC.2}
\]
The right side belongs to
\(k\otimes_{\mathbf F_p}\mathcal O_{\psi,\mathbf F_p}\),
which is naturally \(\mathcal O_{\psi,k}\).
Every nonzero first input on the right is
\(\mathbf F_p\)-valued and cuts out a cyclic extension
of degree \(p\) after restriction to \(K=E(\mu_p)\).
Consequently KT-CUBIC-P-ONLY-NORM expresses every term
by its matched norm-ideal class.  Formula (FC.2)
therefore extends the exact norm comparison to all
finite coefficient fields.
\end{proposition}

\begin{proof}
The characters in (FC.1) take values in \(\mathbf F_p\).
Continuous cochains with finite \(k\)-coefficients
are obtained from their \(\mathbf F_p\)-versions
by scalar extension: a basis of \(k/\mathbf F_p\)
identifies them with a finite direct sum.
Cohomology, all local maps, their kernels, and the
relative pair quotient (CM.14) consequently commute
with that scalar extension.

Choose \(\mathbf F_p\)-linear cocycle representatives
of the two finite spaces in (FC.1).
For every pair \((u_i,v_j)\), choose an ordinary
quadratic defining cochain \(w_{ij}\) as in
KT-GLOBAL-QUADRATIC-MIXED-PERSISTENCE.  Define
\[
                 w(u,v)=\sum_{i,j}c_i d_j w_{ij}.
                                                        \tag{FC.3}
\]
Its differential has the required value because
the quadratic defining equation is bilinear.
The ordinary condition is bilinear as well.
Indeed the local primitive of the strict cocycle
\(v\) at \(p\) can be chosen linearly: it is unique
because \(H^0(\mathbf Q_p,\mathbf F_p(\psi^{-1}))=0\).
The correction \(d(v)(1+\beta)u/2\) in (CM.8)
is then bilinear, and (FC.3) retains the required
unramified class.  Away from \(p\), each chosen
quadratic family is literally trivial on inertia
over \(E\).  Its logarithmic coordinates vanish
there, so (FC.3) retains this condition too.
The uniqueness modulo global coboundary proved
before (CM.8) shows that this choice computes
the same intrinsic ordinary obstruction.

With (FC.3), the cocycle \(-\Omega_+\) in (CM.10)
is homogeneous of degree two in \(u\) and degree
one in \(v\): its terms are \(uw\) and \(u^2v\).
We must retain the same homogeneity in the local
primitives, rather than merely in its raw class.
At \(p\), every primitive is equivalent modulo
the allowed full local \(H^1\), and local \(H^2\)
vanishes.  Choose a linear primitive on the
finite span of the two-cocycles at issue.
At \(q\mid D\), the quadratic family factors
through \(G_{\mathbf Q_q}/I_{E,q}\).
The defect therefore factors through that group.
It has \(p\)-cohomological dimension one, so on
the same finite span choose a linear primitive
within its continuous cochain complex.
Such a primitive is precisely an acceptable
unramified-over-\(E\) cubic local lift in the
rational upper direction.  Its ambiguity is
the allowed unramified local \(H^1\).
Thus the relative pair, not just its raw
two-cocycle, can be chosen homogeneous of
bidegree \((2,1)\) in \((u,v)\).

Passing to the linear relative quotient proves
that \(q\) is quadratic in the first input and
linear in the second, over \(\mathbf F_p\),
with the same operation after scalar extension.
Expanding a quadratic map in a basis gives
exactly (FC.2); the off-diagonal coefficient
is its polarization
\(q(u_i+u_{i'},v_j)-q(u_i,v_j)-q(u_{i'},v_j)\).
No division by a field-dependent scalar is used.

Basis vectors \(u_i\) are nonzero, and
\(u_i+u_{i'}\ne0\) for distinct basis vectors.
Prime-to-\(p\) restriction to \(K\) preserves
their nonzero cohomology classes.  Their
\(\mathbf F_p\)-valued images are therefore
cyclic of order \(p\), exactly the scope of
the matched norm theorem.  If \(U\) or \(V\)
is zero, the formula is the empty sum and
the obstruction is zero.  Applying (CE.2)
and (CE.11) to the finitely many terms proves
the claimed norm comparison for arbitrary
finite \(k\).
\end{proof}

This proposition supplies the scalar reduction
needed for the finite-coefficient norm formula.
It evaluates none of its ideal classes, does
not prove the operation injective in \(v\),
and does not supply an actual source-to-jet map.
